\documentclass[11pt]{article}

\usepackage[margin=1in]{geometry}
\usepackage[T1]{fontenc}
\usepackage[utf8]{inputenc}
\usepackage{amsmath,amssymb,amsthm}
\usepackage{mathtools}
\usepackage{xcolor}
\usepackage{hyperref}
\usepackage{enumitem}
\usepackage{tcolorbox}
\tcbuselibrary{breakable}

\hypersetup{
  colorlinks=true,
  linkcolor=blue!50!black,
  urlcolor=blue!50!black,
  citecolor=blue!50!black
}

\setlength{\parindent}{0pt}
\setlength{\parskip}{0.6em}

% Core notation retained from the source document
\newcommand{\Worlds}{\mathcal{W}}
\newcommand{\Ideas}{\mathcal{I}}
\newcommand{\Noetics}{\mathcal{N}}
\newcommand{\Foundations}{\mathcal{F}}
\newcommand{\Acquisitions}{\mathcal{A}}
\newcommand{\noe}[1]{\nu_{#1}}
\newcommand{\Frac}{\Phi}
\newcommand{\RPM}{\mathfrak{R}}
\newcommand{\Lacunarity}{\Lambda}
\newcommand{\PsychAttractor}{\mathcal{A}_{\psi}}
\newcommand{\ProfileFrac}{\Phi_{\mathrm{profile}}}
\newcommand{\VulnScore}{\mathcal{V}}
\newcommand{\LevelI}{\mathcal{L}_1}
\newcommand{\LevelII}{\mathcal{L}_2}
\newcommand{\LevelIII}{\mathcal{L}_3}
\newcommand{\LevelIV}{\mathcal{L}_4}
\newcommand{\AIGoal}{\mathcal{G}_{\mathrm{AI}}}
\newcommand{\HumanValue}{\mathcal{V}_{\mathrm{H}}}
\newcommand{\Noetica}{\mathbf{Noetica}}
\newcommand{\Sh}{\mathbf{Sh}}

\theoremstyle{plain}
\newtheorem{proposition}{Proposition}

\newtcolorbox{claimbox}[2]{
  breakable,
  colback=gray!5!white,
  colframe=gray!55!black,
  title=\textbf{#1}: #2
}

\title{\textbf{TKS Cognitive Framework}\\[0.5em]
\large Referee-Safe Reduction and Claim Classification}
\author{Prepared from the December 2025 source text}
\date{April 8, 2026}

\begin{document}
\maketitle

\begin{abstract}
This document is a referee-safe reduction of the source file
\texttt{TKS\_FULL\_Cognitive\_Framework.tex}. Its purpose is editorial rather
than promotional: it preserves only those statements that can be honestly
presented as definitions, adopted axioms, imported standard mathematics,
explicit threat models, heuristics, or short proved consequences of stated
definitions. Claims of empirical validity, predictive power, psychological
diagnosis, population-level control, or resolved AI safety are removed or
reclassified as unsupported application hypotheses. The result is not a proof
of the source document's psychological, social, or AI claims. It is a narrower
artifact designed to separate formal structure from interpretation and to mark
clearly where validation is still missing.
\end{abstract}

\section{Scope and Labeling Rule}

\begin{claimbox}{Definition}{Editorial rule}
Every substantive retained claim in this reduction is labeled as exactly one of:
\textbf{Definition}, \textbf{Axiom}, \textbf{Imported Standard Result},
\textbf{Threat Model}, \textbf{Heuristic}, \textbf{Application Hypothesis}, or
\textbf{Proved Statement}.
\end{claimbox}

\begin{claimbox}{Definition}{Referee-safe meaning}
``Referee-safe'' here means only that the claim surface is honest and bounded.
It does \emph{not} mean that the underlying framework has been empirically
validated.
\end{claimbox}

\begin{claimbox}{Threat Model}{Intended use}
Material retained below for psychology, social influence, and predictive
analysis is presented only as a defensive or analytic threat-model language.
Nothing in this document should be read as an endorsed deployment recipe.
\end{claimbox}

\section{Formal Core Retained From the Source}

\begin{claimbox}{Definition}{Noetic alphabet}
Adopt a finite alphabet of ten labeled operators
\[
\Noetics = \{\noe{0},\noe{1},\ldots,\noe{9}\}.
\]
In this reduction, the operators are treated as named symbols inside a formal
model. No empirical semantics are assumed beyond the labels chosen by the
framework.
\end{claimbox}

\begin{claimbox}{Axiom}{D/W/P acquisition order}
For each foundation index $m$, the acquisition predicates satisfy
\[
\sigma(P_m)\Rightarrow \sigma(W_m)\Rightarrow \sigma(D_m).
\]
This is retained as a proposed structural axiom of the framework, not as an
empirical law of human cognition.
\end{claimbox}

\begin{claimbox}{Definition}{Chain level}
For each foundation $m$, define a chain level by
\[
\mathrm{Level}_m =
\begin{cases}
0 & \text{if } \neg \sigma(D_m),\\
1 & \text{if } \sigma(D_m)\land \neg \sigma(W_m),\\
2 & \text{if } \sigma(W_m)\land \neg \sigma(P_m),\\
3 & \text{if } \sigma(P_m).
\end{cases}
\]
\end{claimbox}

\begin{claimbox}{Definition}{Profile sequence}
A profile sequence is a finite or countable word
\[
\ProfileFrac = \langle k_1:k_2:\cdots:k_n\rangle,\qquad k_i\in\{0,\ldots,9\}.
\]
This is a formal encoding device only. No automatic extraction validity is
claimed here.
\end{claimbox}

\begin{claimbox}{Definition}{Grouped domain scores}
If a positive 40-state vector $x=(x_1,\ldots,x_{40})$ satisfies
$\sum_{i=1}^{40}x_i=1$, define four grouped domain scores
\[
d_j(x)=\sum_{i=10(j-1)+1}^{10j}x_i,\qquad j=1,2,3,4.
\]
\end{claimbox}

\begin{proposition}[Grouped simplex property]
If $x_i>0$ for all $i$ and $\sum_i x_i=1$, then each $d_j(x)>0$ and
$\sum_{j=1}^{4} d_j(x)=1$.
\end{proposition}

\begin{proof}
Each $d_j(x)$ is a sum of positive terms, so $d_j(x)>0$. Summing the four block
sums simply partitions the total sum of the forty coordinates, hence
\[
\sum_{j=1}^{4}d_j(x)=\sum_{i=1}^{40}x_i=1.
\]
\end{proof}

\begin{proposition}[Immediate consequence of the D/W/P axiom]
Under the retained acquisition axiom, $\mathrm{Level}_m=3$ implies
$\mathrm{Level}_m\ge 2$, and $\mathrm{Level}_m\ge 2$ implies
$\mathrm{Level}_m\ge 1$.
\end{proposition}

\begin{proof}
If $\mathrm{Level}_m=3$, then $\sigma(P_m)$ holds, hence by the axiom
$\sigma(W_m)$ and $\sigma(D_m)$ also hold. Therefore the chain has at least the
properties required for the lower levels. The second implication is identical.
\end{proof}

\section{Imported Mathematical Templates}

\begin{claimbox}{Imported Standard Result}{Contractive fixed-point principle}
If a self-map on a complete metric space is contractive, then it has a unique
fixed point and iterates converge to it. This is standard mathematics. It does
\emph{not} by itself validate any psychological or social interpretation.
\end{claimbox}

\begin{claimbox}{Definition}{Lacunarity template}
For a suitable box-mass function $M_\varepsilon$, one may define
\[
\Lacunarity(\Frac,\varepsilon)
= \frac{\mathrm{Var}[M_\varepsilon]}{\mathbb{E}[M_\varepsilon]^2}+1.
\]
This is a mathematically standard kind of gap statistic in fractal analysis.
\end{claimbox}

\begin{claimbox}{Imported Standard Result}{Sheaf language}
If a site $(\mathcal{C},J)$ is given, then the category $\Sh(\mathcal{C},J)$ of
sheaves is a Grothendieck topos. This is standard category theory. Any use of
sheaf language for values or cognition requires additional modeling choices.
\end{claimbox}

\begin{claimbox}{Imported Standard Result}{Graph-based threshold dynamics}
Threshold and cascade models on graphs are a standard mathematical template for
studying diffusion processes. Their use in any concrete social or cognitive
application depends on justified state variables, parameters, and data.
\end{claimbox}

\section{Retained Threat Models and Heuristics}

\begin{claimbox}{Heuristic}{NLP-to-noetic mapping}
One may attempt to map text or behavior into profile sequences
$\ProfileFrac$ by a classifier over the ten-symbol alphabet. In this reduction,
that mapping is treated as a heuristic feature-extraction step. No claim is made
that such a mapping is unique, validated, or psychologically complete.
\end{claimbox}

\begin{claimbox}{Threat Model}{Psychological profiling}
Given a chosen extraction pipeline, one can define a workflow
\[
\text{text/behavior} \longrightarrow \ProfileFrac \longrightarrow
\text{derived features}
\]
for defensive analysis of possible profiling systems. This is retained only as a
threat model describing how an adversary might operationalize a symbolic
framework.
\end{claimbox}

\begin{claimbox}{Heuristic}{Attractor language for profiles}
If a profile update map is \emph{assumed} to be contractive on a chosen state
space, then fixed-point language may be used to speak about an attractor-like
summary state. In the absence of a validated state space and validated dynamics,
this remains heuristic rather than diagnostic science.
\end{claimbox}

\begin{claimbox}{Heuristic}{Vulnerability score}
The source document's vulnerability score based on lacunarity is retained only
as a heuristic ranking function. Thresholds such as ``critical if
$\VulnScore>1.5$'' are not preserved as validated scientific cutoffs.
\end{claimbox}

\begin{claimbox}{Threat Model}{Influence and cascade modeling}
Graph diffusion, seed-node choice, campaign propagation, and attractor
perturbation can be written down as formal threat models for how a hostile actor
might reason about influence operations. They are not presented here as proven
laws of populations.
\end{claimbox}

\begin{claimbox}{Threat Model}{Predictive pipeline}
The six-stage pipeline from digital traces to inferred goals, trajectories, and
intervention points is retained only as an adversarial capability model. It is
not retained as an established predictive system.
\end{claimbox}

\section{Application Hypotheses for AI}

\begin{claimbox}{Application Hypothesis}{Alignment decomposition}
The source proposes that alignment analysis may be \emph{modeled} by decomposing
goals into desire, wisdom, and power components and expressing constraints
through an RPM-like wrapper. This is a design hypothesis, not a solved
alignment theorem.
\end{claimbox}

\begin{claimbox}{Application Hypothesis}{Interpretability vocabulary}
The ten-symbol noetic alphabet may be explored as a human-readable vocabulary
for annotating model behavior. This is a proposed interpretive schema, not an
established decomposition of neural computation.
\end{claimbox}

\begin{claimbox}{Application Hypothesis}{Value-sheaf viewpoint}
It may be mathematically convenient to describe context-dependent value
assignments using sheaf-like language over a chosen site. This is a modeling
proposal only; no result here shows that real value learning in AI is solved by
such a construction.
\end{claimbox}

\begin{claimbox}{Application Hypothesis}{Safety architecture}
Temporal typing, coalgebraic dynamics, or categorical organization may be useful
for specifying AI safety architectures. The source document does not provide the
formal verification machinery or empirical evidence needed to claim deployment
readiness.
\end{claimbox}

\section{Claims Explicitly Removed From the Referee-Safe Version}

\begin{claimbox}{Definition}{Removed claim class}
The following kinds of claims from the source document are intentionally omitted
from the referee-safe version because they were unsupported, overstated, or too
operational for a scholarly artifact.
\end{claimbox}

\begin{itemize}[leftmargin=1.5em]
\item theorem-style claims without proofs, including the ``40-Step
Stabilization'' claim and application-level cascade or predictive theorems;
\item assertions that TKS provides ``rigorous solutions'' to core AI crises;
\item statements that collective attractors, value resolution, or uncertainty
tracking make AI systems robust or safe ``by construction'';
\item operational scenario language that reads like a deployment recipe for
psychological manipulation, population steering, or mass surveillance;
\item threshold claims presented as if empirically calibrated when no evidence
is supplied.
\end{itemize}

\section{What Would Be Required for Validation}

\begin{claimbox}{Definition}{Validation requirements}
To move any retained heuristic or application hypothesis toward scientific
credibility, the framework would need at minimum:
\begin{enumerate}[label=(\arabic*)]
  \item an explicit state space and observation model;
  \item reproducible extraction procedures from data to symbolic states;
  \item benchmark datasets and inter-rater agreement where human annotation is
  involved;
  \item predictive evaluation against held-out outcomes;
  \item ablations against simpler baselines;
  \item false-positive and false-negative analysis for any risk or vulnerability
  scoring.
\end{enumerate}
\end{claimbox}

\section{Conclusion}

\begin{claimbox}{Definition}{Honest status}
This referee-safe reduction does \emph{not} certify the original cognitive
framework as scientifically valid. It preserves a narrower core: formal
definitions, a few trivial proved consequences, standard imported mathematics,
and clearly labeled heuristics or threat models. Everything stronger remains
future work or unsupported interpretation.
\end{claimbox}

\begin{thebibliography}{9}
\bibitem{barnsley}
M. F. Barnsley, \emph{Fractals Everywhere}, Academic Press, 1988.

\bibitem{banach}
S. Banach, ``Sur les opérations dans les ensembles abstraits et leur application
aux équations intégrales,'' \emph{Fundamenta Mathematicae} 3 (1922), 133--181.

\bibitem{maclane}
S. Mac Lane and I. Moerdijk, \emph{Sheaves in Geometry and Logic}, Springer,
1992.

\bibitem{threshold}
D. Centola and M. Macy, ``Complex contagions and the weakness of long ties,''
\emph{American Journal of Sociology} 113(3) (2007), 702--734.
\end{thebibliography}

\end{document}
